% Generated draft: requires mathematical and source review.
\documentclass[11pt]{article}
\usepackage[margin=25mm]{geometry}
\usepackage{fontspec}
\setmainfont{Latin Modern Roman}
\usepackage{amsmath,amssymb,mathtools,unicode-math}
\setmathfont{Latin Modern Math}
\usepackage{enumitem,needspace,xurl,hyperref}
\setlength{\parindent}{0pt}
\setlength{\parskip}{5pt}
\setlength{\emergencystretch}{3em}
\newcommand{\N}{\mathbb{N}}
\newcommand{\Z}{\mathbb{Z}}
\newcommand{\Q}{\mathbb{Q}}
\newcommand{\R}{\mathbb{R}}
\newcommand{\C}{\mathbb{C}}
\newcommand{\F}{\mathbb{F}}
\newcommand{\A}{\mathbb{A}}
\newcommand{\PP}{\mathbb{P}}
\newcommand{\E}{\mathbb{E}}
\newcommand{\eps}{\varepsilon}
\newcommand{\dd}{\,\mathrm{d}}
\begin{document}
\section*{UnsolvedMath ID: 10400061}
\textbf{Source problem code:} catalogue entry
\subsection{Definitions and mathematical statement}
Configuration-space integral constructions use compactified configuration spaces with boundary strata describing collisions of points. In a standard terminology, hidden collision strata involve larger clusters rather than only the principal pairwise faces; the precise strata in the three approaches referenced by the source are not supplied.

Find another way to eliminate the hidden-stratum contributions so that all three source approaches naturally become the mapping degree of one and the same map. For context, a degree interpretation for a map $F:X\to Y$ between closed oriented spaces of equal dimension satisfies $F_*[X]=\deg(F)[Y]$; configuration-space versions may require gluing, relative classes, or other specified conventions. The actual spaces, target, map, and boundary cancellation data are absent. The aim is a common geometric degree construction, not merely numerical equality of three unnamed invariants.

The extract records two qualified pieces of background: S. Poirier [328] related the first and second approaches using quasi-tangles under a vanishing-anomaly assumption, with a cross-reference to Question 3.12; D. Thurston suggested Etingof--Kazhdan $R$-matrices [117] as potentially useful for comparing the first and third. The former is conditional on the stated anomaly hypothesis and the latter is a suggestion, not an asserted equivalence. Neither is promoted into the requested alternative hidden-stratum construction. The following Section 3.6 heading is separate material.
\subsection{Short English statement}
Eliminate hidden configuration-space boundary contributions in a new way that unifies the source's three constructions as the degree of the same map. Their definitions and the exact boundary/map data are missing; the stated Poirier comparison remains conditional on vanishing anomaly.
\subsection{Sources}
\begin{itemize}
\item Ohtsuki, \emph{Problems on invariants of knots and 3-manifolds} (2002), Problem 3.11, PDF page 56, attributed to M. Polyak, \url{https://msp.org/gtm/2002/04/gtm-2002-04-024s.pdf}; accessed 2026-07-29, PDF-font-structure extraction.
\item S. Poirier [328], conditional first/second-approach comparison via quasi-tangles, and source Question 3.12; D. Thurston's suggestion involving Etingof--Kazhdan $R$-matrices [117]. The record supplies no bibliography entries for [328] or [117].
\item Local registry \texttt{data/unsolvedmath/problems.json}, permanent UnsolvedMath ID $10400061$, source code catalogue entry, \emph{Problem 3.11}. Dataset attribution: Ulam AI and the UnsolvedMath Contributors.
\end{itemize}
\end{document}
